\section{Verifier-hacking audit}
\label{app:hacking}

If GRPO were exploiting the reward checker rather than solving the task, its
accuracy would be an artifact and the comparison meaningless. We audit the
reward-dump records of the 14 GRPO runs whose dumps were retained, spanning all
three backbones. H1 flags a constant program that performs no computation; H2
flags the gold value appearing in the emitted code as a literal; H3 is the
share of positive samples falling under the single most frequent code skeleton;
H4 flags no-op arithmetic such as $x \times 1$ or $x + 0$, which would slip past
the anti-constant guard.

\begin{table}[t]\centering\small
\setlength{\tabcolsep}{3.5pt}
\caption{Verifier-hacking audit over the 14 GRPO runs whose reward dumps were
retained. H1: constant program; H2: gold value appearing as a literal; H3:
share of positive samples under the single most frequent code skeleton; H4:
no-op arithmetic. Counts are over positive-reward samples.}
\label{tab:supp-hacking}
\begin{tabular}{c rrr rrr r}
\toprule
cfg & runs & samples & positive & H1 & H2 & H4 & H3 \\
\midrule
1 & 6 & 87{,}872 & 47{,}616 & 0 & 90 & 4 & $\leq 0.9\%$ \\
2 & 6 & 37{,}933 & 6{,}030 & 0 & 14 & 23 & $\leq 1.5\%$ \\
3 & 2 & 4{,}672 & 3{,}871 & 0 & 7 & 0 & $\leq 2.1\%$ \\
\midrule
total & 14 & 130{,}477 & 57{,}517 & \textbf{0} & 111 & 27 & $\leq 2.1\%$ \\
\bottomrule
\end{tabular}
\end{table}


\noindent
Constant-program shortcuts occur \textbf{zero} times in 57{,}517
positive-reward samples: the verifier's anti-constant guard holds. No-op
arithmetic appears 27 times, $0.05\%$ of positive samples. Template collapse is
absent---the most frequent skeleton accounts for at most $2.1\%$ of positive
samples, with skeleton Gini coefficients between $0.06$ and $0.45$, indicating
a broadly distributed solution space rather than a single exploited pattern.

H2 fires on $0.19\%$ of positive samples. This is expected rather than
suspicious on these tasks: gold answers are computed from numbers that appear
in the context, so a correct program will sometimes contain a literal that
coincides with the gold value.

The four unretained dumps are Config~3 runs. Their positive-reward rates and
group statistics are nonetheless available from the run metadata and appear in
Table~\ref{tab:zeroadv}, so no run is unaccounted for. The audit therefore
covers 14 of 18 runs by dump inspection and 18 of 18 by recorded aggregate.
